\documentclass[../calculus.tex]{subfiles}

\begin{document}


\chapter{Limits and Continuity}


\section{Limits of Functions}

The fundamental concept of calculus is the limit. We formalise the notion of a function approaching a specific value. All of this content is covered in depth in Analysis I (MATH40002) so we will not spend much time on it, and all proofs can be found in Analysis I.

\begin{definition}[$\varepsilon$-$\delta$ definition of a limit]
    \begin{shaded*}
    Let $f$ be a function defined on an open interval containing $x_0$, except possibly at $x_0$ itself. We say that $\lim_{x\to x_0}f(x)=\ell$ if for every $\varepsilon>0$, there exists a $\delta>0$ such that for all $x$, 
    \[0<|x-x_0|<\delta \implies |f(x)-\ell|<\varepsilon.\]
    \end{shaded*}
\end{definition}

\begin{example}
    We prove that $\lim_{x\to 2}(3x-1)=5$. Given $\varepsilon>0$, we want to find $\delta>0$ such that $0<|x-2|<\delta$ implies $|(3x-1)-5|<\varepsilon$. We manipulate the absolute value we are trying to bound:
    $$|(3x-1)-5| = |3x-6| = 3|x-2|.$$
    To ensure this is strictly less than $\varepsilon$, we require $3|x-2|<\varepsilon$, which is equivalent to $|x-2|<\frac{\varepsilon}{3}$. Therefore, we choose $\delta = \frac{\varepsilon}{3}$. Then for all $x$ satisfying $0<|x-2|<\delta$, we have $|(3x-1)-5| = 3|x-2| < 3\delta = \varepsilon$.
\end{example}

\begin{theorem}[Algebra of Limits]
    \begin{shaded*}
        Suppose $\lim_{x\to x_0}f(x)=L$ and $\lim_{x\to x_0}g(x)=M$. Then:
        \begin{enumerate}
            \item $\lim_{x\to x_0} (f(x)+g(x)) = L+M$
            \item $\lim_{x\to x_0} (f(x)g(x)) = LM$
            \item $\lim_{x\to x_0} \frac{f(x)}{g(x)} = \frac{L}{M}$, provided $M \neq 0$
            \item $\lim_{x\to x_0} (cf(x)) = cL$ for any constant $c\in\R$
        \end{enumerate}
    \end{shaded*}
\end{theorem}

\begin{definition}[Limit at infinity]
    \begin{shaded*}
    We say that $\lim_{x\to\infty}f(x)=\ell$ if for every $\varepsilon>0$, there exists an $A>0$ such that for all $x$,
    \[x>A \implies |f(x)-\ell|<\varepsilon.\]
    \end{shaded*}
\end{definition}
One makes the analogous definition is made for limits at $-\infty$.

\begin{example}
    We prove $\lim_{x\to\infty}\frac{1}{x}=0$. Given $\varepsilon>0$, we require $\left|\frac{1}{x}-0\right|<\varepsilon$, which means $\frac{1}{x}<\varepsilon$ for $x>0$. Rearranging gives $x>\frac{1}{\varepsilon}$. We choose $A=\frac{1}{\varepsilon}$. For all $x>A$, we have $x>\frac{1}{\varepsilon}$, so $\frac{1}{x}<\varepsilon$.
\end{example}

\begin{definition}
    \begin{shaded*}
    We say that $\lim_{x\to x_0}f(x)=\infty$ if for every $B>0$, there exists a $\delta>0$ such that for all $x$,
    \[0<|x-x_0|<\delta \implies f(x)>B.\]
    \end{shaded*}
\end{definition}
Again, one makes the analogous definition for when $\lim_{x\to x_0}f(x)=-\infty$, and when $x\to \pm \infty$.


\begin{definition}[One-sided limits]
    \begin{shaded*}
    The right-sided limit $\lim_{x\to x_0^+}f(x)=\ell$ if for every $\varepsilon>0$, there exists a $\delta>0$ such that for all $x$,
    \[0<x-x_0<\delta \implies |f(x)-\ell|<\varepsilon.\]
    The left-sided limit $\lim_{x\to x_0^-}f(x)=\ell$ if for every $\varepsilon>0$, there exists a $\delta>0$ such that for all $x$,
    \[0<x_0-x<\delta \implies |f(x)-\ell|<\varepsilon.\]
    A two-sided limit exists if and only if both one-sided limits exist and are equal.
    \end{shaded*}
\end{definition}
\begin{example}
    Let $H:\R\to\R$ (the so-called Heaviside function) be defined by \[H(x)=\begin{cases}
        1, & x\ge 0,\\
        0, & x<0.
    \end{cases}\]
    We see that $\lim_{x\to 0^-}H(x)=0$ and $\lim_{x\to 0^+}H(x)=1$, and $H(0)=1$. Even though $H(0)=1$, the limit approaching 0 from the left is still zero.
\end{example}

\begin{theorem}[Squeeze Theorem]
    \begin{shaded*}
        Suppose $g(x) \le f(x) \le h(x)$ for all $x$ in an open interval containing $x_0$, except possibly at $x_0$ itself. If $\lim_{x\to x_0}g(x) = \lim_{x\to x_0}h(x) = \ell$, then $\lim_{x\to x_0}f(x) = \ell$.
    \end{shaded*}
\end{theorem}

\begin{example}
    Evaluate $\lim_{x\to 0} x \sin\left(\frac{1}{x}\right)$. 
    Since $-1 \le \sin\left(\frac{1}{x}\right) \le 1$ for all $x \neq 0$, we multiply the inequality by $|x|$ to obtain:
    \[-|x| \le x \sin\left(\frac{1}{x}\right) \le |x|.\]
    We know $\lim_{x\to 0} (-|x|) = 0$ and $\lim_{x\to 0} |x| = 0$. By the comparison test, $\lim_{x\to 0} x \sin\left(\frac{1}{x}\right) = 0$.
\end{example}

\begin{example}

A famous limit that is an indeterminate form but can be derived using geometric ideas is the following:
\[\lim_{x\to 0}\frac{\sin x}{x}.\]

Consider a unit circle centered at the origin $O(0,0)$. Let $x$ be a small positive angle such that $0 < x < \frac{\pi}{2}$, measured in radians. Let $A$ be the point $(1,0)$ and $B$ be the point on the circle such that $\angle AOB = x$. The coordinates of $B$ are $(\cos x, \sin x)$.

Draw a vertical tangent line to the circle at point $A$. Extend the ray $OB$ until it intersects this tangent line at point $C$. The coordinates of $C$ are $(1, \tan x)$.

We compare the areas of three geometric regions:
\begin{enumerate}
    \item The triangle $\triangle OAB$, which has base $1$ and height $\sin x$.
    \item The circular sector $OAB$, which has radius $1$ and central angle $x$.
    \item The triangle $\triangle OAC$, which has base $1$ and height $\tan x$.
\end{enumerate}

From the geometric construction, the triangle $\triangle OAB$ is contained within the sector $OAB$, which in turn is contained within the triangle $\triangle OAC$. Therefore, their areas satisfy the inequality:
$$
\text{Area}(\triangle OAB) < \text{Area}(\text{sector } OAB) < \text{Area}(\triangle OAC).
$$
Substituting the formulas for these areas yields:
$$
\frac{1}{2}(1)(\sin x) < \frac{1}{2}(1)^2 x < \frac{1}{2}(1)(\tan x).
$$
Multiplying by $2$, we obtain:
$$
\sin x < x < \tan x.
$$
Since $0 < x < \frac{\pi}{2}$, we know $\sin x > 0$. Dividing the entire inequality by $\sin x$ gives:
$$
1 < \frac{x}{\sin x} < \frac{1}{\cos x}.
$$
Taking the reciprocal of all terms reverses the inequalities:
$$
\cos x < \frac{\sin x}{x} < 1.
$$
As $x \to 0^+$, we know $\lim_{x \to 0^+} \cos x = 1$ and $\lim_{x \to 0^+} 1 = 1$. By the Squeeze Theorem, it follows that:
$$
\lim_{x \to 0^+} \frac{\sin x}{x} = 1.
$$
Since $f(x) = \frac{\sin x}{x}$ is an even function $\left(\text{as } \frac{\sin(-x)}{-x} = \frac{-\sin x}{-x} = \frac{\sin x}{x}\right)$, the left-hand limit is identical to the right-hand limit:
$$
\lim_{x \to 0^-} \frac{\sin x}{x} = 1.
$$
Therefore, the two-sided limit evaluates to $1$:
$$
\lim_{x \to 0} \frac{\sin x}{x} = 1. 
$$
\end{example}

\begin{example}
Another simple trigonometric limit is \[\lim_{x\to 0}\frac{\cos x-1}{x},\]
which again is of an indeterminate form, and so care must be taken to find this limit.

We evaluate the limit by multiplying the numerator and denominator by the conjugate of the numerator, $\cos x + 1$:
\begin{align*}
\lim_{x \to 0} \frac{\cos x - 1}{x} &= \lim_{x \to 0} \frac{(\cos x - 1)(\cos x + 1)}{x(\cos x + 1)} \\
&= \lim_{x \to 0} \frac{\cos^2 x - 1}{x(\cos x + 1)}.
\end{align*}
Using the Pythagorean identity $\sin^2 x + \cos^2 x = 1$, we substitute $\cos^2 x - 1 = -\sin^2 x$:
\begin{align*}
\lim_{x \to 0} \frac{-\sin^2 x}{x(\cos x + 1)} &= \lim_{x \to 0} \left( -\frac{\sin x}{x} \cdot \frac{\sin x}{\cos x + 1} \right).
\end{align*}
Using the algebra of limits and the result from the previous proof:
\begin{align*}
\lim_{x \to 0} \left( -\frac{\sin x}{x} \cdot \frac{\sin x}{\cos x + 1} \right) &= -\left( \lim_{x \to 0} \frac{\sin x}{x} \right) \cdot \left( \lim_{x \to 0} \frac{\sin x}{\cos x + 1} \right) \\
&= -(1) \cdot \left( \frac{0}{1 + 1} \right) \\
&= 0.
\end{align*}
Therefore, $\lim_{x \to 0} \frac{\cos x - 1}{x} = 0$.

\end{example}

\section{Continuity}

\begin{definition}[Continuity at a point]
    \begin{shaded*}
    A function $f$ is continuous at $x_0$ if for every $\varepsilon>0$, there exists a $\delta>0$ such that for all $x$ in the domain of $f$,
    \[|x-x_0|<\delta \implies |f(x)-f(x_0)|<\varepsilon.\]
    Equivalently, this requires that $\lim_{x\to x_0}f(x)=f(x_0)$.
    \end{shaded*}
\end{definition}

\begin{remark}
    For $f$ to be continuous at $x_0$, three conditions must be satisfied: $f(x_0)$ must be defined, $\lim_{x\to x_0}f(x)$ must exist, and these two values must be equal.
\end{remark}

\begin{example}
    Let $f(x)$ be defined as:
    $$f(x) = \begin{cases} 
    \frac{x^2-1}{x-1} & x \neq 1 \\
    2 & x = 1 
    \end{cases}$$
    We check for continuity at $x=1$. We have $f(1) = 2$. Taking the limit as $x \to 1$, we note that $x \neq 1$, so we can divide by $x-1$:
    \[\lim_{x\to 1} \frac{x^2-1}{x-1} = \lim_{x\to 1} \frac{(x-1)(x+1)}{x-1} = \lim_{x\to 1} (x+1) = 2.\]
    Since $\lim_{x\to 1} f(x) = f(1)$, the function is continuous at $x=1$.
\end{example}
There is an analogous version of the Algebra of Limits for continuous functions (as one would expect, since continuity is defined in terms of limits):
\begin{theorem}[Algebra of Continuous Functions]
    \begin{shaded*}
        Let $f,g$ be continuous functions. Then
        \begin{enumerate}
            \item[1.] $f+g$ is continuous.
            \item[2.] $fg$ is continuous.
            \item[3.] $\frac{f}{g}$ is continuous, provided $g(x)\ne 0$ for all $x$.
            \item[4.] $g\circ f$ is continuous.   
        \end{enumerate}
    \end{shaded*}
\end{theorem}
There are pathological examples of functions which are discontinuous everywhere on $\R$. For example, consider the function $f:\R\to\R$ defined by
\[f(x)=\begin{cases}
    1,& x\in\Q,\\
    0,& x\notin\Q.
\end{cases}\]
By density of the rational and irrational numbers in $\R$, one checks that $f$ is discontinuous everywhere (exercise).




\end{document}